% How this book was made, and the table of the twenty contributions.
\chapter{How this book was made}
\markboth{How this book was made}{How this book was made}

\section{Who wrote what}

This book was written by AI agents (Claude, by Anthropic) under the direction of the author. One agent designed the
frame: the layout, the prologue, the chapter on the method, the synthesis chapter, the letter and these closing pages.
Four other agents wrote the twenty contribution chapters in parallel, from the same style guide and the same sheet for
each contribution, and with one strict rule: use no fact and no number that does not appear in the records of the
work itself. The author chose the subject, the tone and the dedication, and he answers for the content. At the time
of writing, no outside scientist has reviewed this book.

This English edition is a translation of the French edition, \textit{Le journal des découvertes}, also made by AI
agents under the author's direction: one translated the frame, and four others translated the contribution chapters
in parallel. They tried to keep the same
claims, the same levels and the same limits as the French text, but some wording may still differ from it.

The results told here were produced in the same way, by AI agents under the direction of the author, and audited by
other AI agents, following the method described in Chapter~\ref{ch:elenchus}. \textbf{No human mathematician has
reviewed them.}

\section{How the twenty were chosen}

\paragraph{The inventory} Every record published on Zenodo during the year was listed: 144 records, which correspond to
112 distinct pieces of work once duplicates (the same text deposited up to thirteen times) and successive versions of
the same work are merged. The author's GitHub account holds 420 repositories, of which 317 are copies of other
people's projects; 82 are the author's own projects, created or modified during the year (from the end of October 2025
to the end of October 2026), and 57 of those touch on science. A GitHub repository is kept as a candidate only if no
piece of work published on Zenodo already describes it. In all, 113 candidates.

\paragraph{The grid, fixed before scoring} The selection grid was written and frozen on 26 October 2026, \emph{before}
any candidate was scored, and it was not changed afterwards. It scores every candidate from 0 to 3 on seven criteria,
the same grid for all, whatever the field: the rigour achieved (the highest Elenchus level), independent checking,
reproducibility, synthesis between mathematics, physics or biology, calculation and formal proof, honesty (negative
results, published corrections), novelty, and value for non-specialists. Every score is justified in one line, with
the file it rests on.

\paragraph{The fixed rules} Four rules, applied in this order, complete the grid. Pieces of work whose central claim was
withdrawn or refuted are not scored as successes: they are attached to the work that publishes the withdrawal, and
scored as lessons. At least two places are reserved for each of the five fields (mathematics, physics, genetics and
biology, numerical computing and software, AI method and verification). At least three of the twenty must be lessons.
No field may have more than seven places. The fields are those of the inventory: cosmology and teaching are counted
there separately from physics, which explains why physics in the broad sense has nine chapters (seven in physics and
two in cosmology).

\paragraph{What the rules changed} Compared with the twenty best raw scores, the rules brought in four contributions
(the rusty-SUNDIALS instrument, the GWAYA gate, the AutoevolveAI workbench and the stellarator lesson) to meet the
minimums for the ``numerical computing'' and ``method'' fields, and took out four: two because of the cap of seven in
physics, two displaced by the reserved places. Seventeen of the twenty contributions date from 8 to 26 October 2026:
the grid rewards replay, independent audits with published corrections and negative results, which the project only
produced systematically from October onwards. The earlier months come into the book mostly through the lessons. The
complete selection file, with every score and its justification, is public.\reference{Folder
\nolinkurl{results/discovery_journal_2026/}: \texttt{RUBRIC.md}, \nolinkurl{SELECTION_final.md},
\nolinkurl{TOP20_final.json}, \nolinkurl{scores_final.json}, \texttt{RESCORE.md}.}

\paragraph{A score revised along the way} The first scoring relied on the replay register as it stood halfway through
(71 modules replayed). When the replay was finished (93 modules, all passed), the scores that cited the newly replayed
modules were recalculated, with the same grid, without changing it. Only one contribution on the list changed: the
certificate on screened forces dropped out, tied on score, separated by date and then excluded by the cap of seven in
physics, and the numerical refutation of the conjecture QU-1 (Chapter~\ref{ch:c12}) came in. The order of a few
others moved. This book follows the final list; the old and the new are both kept in the folder.

\section{The twenty contributions}

The table gives, for each contribution, its rank, the title of its chapter, its field, its rigour score in the grid
and its DOI. The rigour score measures the best-checked claim that the work itself states: 3 for a proof checked by
the judge and replayed, or for a calculation certificate rechecked by a different program; 2 for a proof checked by
the judge with an audit of the axioms, without replay; 1 for an argument, a citation or an exact calculation without
the judge; 0 for decimal numbers only. It is not the level of the work as a whole, and it is not quite an Elenchus
level: the level badges, piece by piece, are in each chapter, and they are what count. (The score of contribution
No.\,20 follows the inventory, which does not cite the list of axioms; the repository's check file gives it, hence
the chapter's As badge.)

\begingroup
\small
\renewcommand{\arraystretch}{1.22}
% generated by code/make_backmatter.py from TOP20_final.json and ORDER.json -- do not edit
\begin{tabularx}{\linewidth}{@{}r>{\raggedright\arraybackslash}X>{\raggedright\arraybackslash}p{2.3cm}c@{}}
\toprule
\textsf{No.} & \textsf{contribution (chapter title), DOI or repository} & \textsf{field} & \textsf{rigour}\\
\midrule
1 & Cosmic distances with an error margin checked by a machine \newline {\scriptsize\color{encre!60}2026-10-08, p.~\pageref{ch:c01}} \newline{\footnotesize \doi{10.5281/zenodo.23247730}} & cosmology & {\sffamily\small 3\,/\,3}\\
2 & A famous tension in cosmology, written as a verified inequality \newline {\scriptsize\color{encre!60}2026-10-10, p.~\pageref{ch:c02}} \newline{\footnotesize \doi{10.5281/zenodo.23284020}} & cosmology & {\sffamily\small 3\,/\,3}\\
3 & The triangle, the cheapest arrangement: verified consequences and a simulation \newline {\scriptsize\color{encre!60}2026-10-11, p.~\pageref{ch:c03}} \newline{\footnotesize \doi{10.5281/zenodo.23293429}} & physics & {\sffamily\small 3\,/\,3}\\
4 & Screened repulsion joins the family of the theorem, and an error is repaired \newline {\scriptsize\color{encre!60}2026-10-25, p.~\pageref{ch:c04}} \newline{\footnotesize \doi{10.5281/zenodo.23646105}} & physics & {\sffamily\small 3\,/\,3}\\
5 & Redoing someone else's calculation, without declaring it right \newline {\scriptsize\color{encre!60}2026-10-26, p.~\pageref{ch:c13}} \newline{\footnotesize \doi{10.5281/zenodo.23653827}} & mathematics & {\sffamily\small 3\,/\,3}\\
6 & A ruler proved straight shows that our own sketch was wrong \newline {\scriptsize\color{encre!60}2026-10-26, p.~\pageref{ch:c05}} \newline{\footnotesize \doi{10.5281/zenodo.23653846}} & physics & {\sffamily\small 3\,/\,3}\\
7 & What chance does a mutation have of taking over a population? \newline {\scriptsize\color{encre!60}2026-10-26, p.~\pageref{ch:c06}} \newline{\footnotesize \doi{10.5281/zenodo.23653850}} & genetics, biology & {\sffamily\small 3\,/\,3}\\
8 & When the seating plan changes nothing about selection \newline {\scriptsize\color{encre!60}2026-10-26, p.~\pageref{ch:c07}} \newline{\footnotesize \doi{10.5281/zenodo.23660071}} & genetics, biology & {\sffamily\small 3\,/\,3}\\
9 & A textbook that follows a theorem all the way to the laboratory \newline {\scriptsize\color{encre!60}2026-10-11, p.~\pageref{ch:c08}} \newline{\footnotesize \doi{10.5281/zenodo.23295142}} & teaching & {\sffamily\small 3\,/\,3}\\
10 & Electron--hole pairs that arrange themselves in triangles, on paper \newline {\scriptsize\color{encre!60}2026-10-25, p.~\pageref{ch:c09}} \newline{\footnotesize \doi{10.5281/zenodo.23646103}} & physics & {\sffamily\small 3\,/\,3}\\
\bottomrule
\end{tabularx}

\bigskip

\begin{tabularx}{\linewidth}{@{}r>{\raggedright\arraybackslash}X>{\raggedright\arraybackslash}p{2.3cm}c@{}}
\toprule
\textsf{No.} & \textsf{contribution (chapter title), DOI or repository} & \textsf{field} & \textsf{rigour}\\
\midrule
11 & The force in an atomically thin semiconductor, as a pile of sand \newline {\scriptsize\color{encre!60}2026-10-25, p.~\pageref{ch:c10}} \newline{\footnotesize \doi{10.5281/zenodo.23646107}} & physics & {\sffamily\small 3\,/\,3}\\
12 & Testing before proving: a natural idea that does not hold up \emph{(lesson)} \newline {\scriptsize\color{encre!60}2026-10-26, p.~\pageref{ch:c12}} \newline{\footnotesize \doi{10.5281/zenodo.23653833}} & physics & {\sffamily\small 3\,/\,3}\\
13 & Perfect at rest, not when it trembles \emph{(lesson)} \newline {\scriptsize\color{encre!60}2026-10-26, p.~\pageref{ch:c11}} \newline{\footnotesize \doi{10.5281/zenodo.23655392}} & physics & {\sffamily\small 3\,/\,3}\\
14 & Four floors out of twelve, and the price of the rest \newline {\scriptsize\color{encre!60}2026-10-26, p.~\pageref{ch:c14}} \newline{\footnotesize \doi{10.5281/zenodo.23656804}} & mathematics & {\sffamily\small 3\,/\,3}\\
15 & Rebuilding someone else's machine in your own workshop \newline {\scriptsize\color{encre!60}2026-10-26, p.~\pageref{ch:c15}} \newline{\footnotesize \doi{10.5281/zenodo.23661381}} & mathematics & {\sffamily\small 3\,/\,3}\\
16 & One general guarantee, three precise certificates \newline {\scriptsize\color{encre!60}2026-10-10, p.~\pageref{ch:c16}} \newline{\footnotesize \doi{10.5281/zenodo.23285631}} & mathematics & {\sffamily\small 3\,/\,3}\\
17 & The calculation engine, and its inspection report \newline {\scriptsize\color{encre!60}2026-05-11, p.~\pageref{ch:c17}} \newline{\footnotesize repository \texttt{rusty-SUNDIALS}} & numerical computing & {\sffamily\small 0\,/\,3}\\
18 & A door that stays closed when in doubt \newline {\scriptsize\color{encre!60}2026-10-03, p.~\pageref{ch:c18}} \newline{\footnotesize \doi{10.5281/zenodo.23119953}} & method, verification & {\sffamily\small 2\,/\,3}\\
19 & The test bench that inspected itself \newline {\scriptsize\color{encre!60}2026-09-19, p.~\pageref{ch:c19}} \newline{\footnotesize repository \texttt{AutoevolveAI}} & method, verification & {\sffamily\small 1\,/\,3}\\
20 & The thermostat was accused of overheating \emph{(lesson)} \newline {\scriptsize\color{encre!60}2026-10-21, p.~\pageref{ch:c20}} \newline{\footnotesize \doi{10.5281/zenodo.23532458}} & numerical computing & {\sffamily\small 1\,/\,3}\\
\bottomrule
\end{tabularx}

\endgroup
